
\section{与有理同伦论的相容性}
\label{sec:7-5}

\subsection{结果陈述}
\label{subsec:7-5-1}

本节旨在证明，由
$\overline{\mathfrak m}_k=\mathfrak m_{k,\beta_0}$
得到的 $A_\infty$ 代数同德拉姆复形同伦等价。
也就是说，我们将证明如下结果。

\begin{theorem}{7.5.1}\label{thm:7-5-1}
设 $\Omega(L)$ 是拉格朗日子流形 $L$ 的德拉姆复形。
则定理~\ref{thm:3-1-9} 中的 $A_\infty$ 代数
在 $\mathbb R$ 上同 $\Omega(L)$ 同伦等价。
此同伦等价在同伦意义下是典范的。
\end{theorem}

用 $[\text{有理同伦}]$ 表示以 CW 复形的有理同伦型为对象的范畴，
用 $[\mathrm{CDGA};\text{同伦}]$ 表示如下范畴：
其对象是交换微分分次代数，
其态射是交换微分分次代数同态的同伦类
（例如见 \textbf{[GrMo81]}）。
把流形对应到它的德拉姆复形，便得到函子
\[
[\text{有理同伦}]
\longrightarrow
[\mathrm{CDGA};\text{同伦}].
\]
D. Quillen \textbf{[Qui67, Qui69]} 与沙利文 \textbf{[Sul78]}
证明了：把这个函子限制在幂零的（例如单连通的）
CW 复形所成的全子范畴上时，
它给出该全子范畴同其像之间的范畴等价。

下面定义范畴 $[\text{相对辛}]$。
它的一个对象由辛流形 $M$、
其中一个可定向拉格朗日子流形 $L$，
以及 $L$ 的一个相对 spin 结构组成。
从 $(M,L)$ 到 $(M',L')$ 的一个态射，
是满足 $f(L)=L'$ 且保持相对 spin 结构的辛微分同胚
$f\colon M\to M'$ 的辛同痕类。
这里，若存在一族辛微分同胚 $f_t$，使得
$f_0=f$、$f_1=f'$，并且对每个 $t$ 都有
$f_t(L)=L'$，就称 $f$ 与 $f'$ 辛同痕。
于是可以定义函子
\[
[\text{相对辛}]
\longrightarrow
[\text{有理同伦}],
\qquad
(M,L)\longmapsto[L],
\]
其中 $[L]$ 表示 $L$ 的有理同伦型。

还要指出，对应
\[
(C,d,\mathbin{\cdot})
\longmapsto
(C,\overline{\mathfrak m}),
\]
其中 $\overline{\mathfrak m}$ 由
式~\eqref{eq:3-2-5} 给出
\translatornote{源书此处回指式 (3.2.13)；
但该式组只定义滤过 $A_\infty$ 结构的模正能量约化，
而把微分分次代数视为 $A_\infty$ 代数的两个非零运算及其符号
明确给在式 (3.2.5)。故译文改正回指。见疑误 A-0611。}，
诱导函子
\[
[\mathrm{DGA};\text{同伦}]
\longrightarrow
[A_\infty;\text{同伦}].
\]
范畴 $[A_\infty;\text{同伦}]$ 的对象是无滤过
$A_\infty$ 代数，态射是 $A_\infty$ 同态的同伦类；
它也可以像第~\ref{subsec:4-3-1} 小节那样记作
$\mathfrak{HA}_\infty$。

事实上，上述函子给出范畴等价。
\textbf{[Fuk02II]} 的推论 9.4（另见 \textbf{[Lef03]}）证明，
任意 $A_\infty$ 范畴都同伦等价于一个微分分次范畴。
其代数版本由范畴版本推出，并说明任意 $A_\infty$ 代数
作为 $A_\infty$ 代数都同伦等价于一个微分分次代数。
这也可由如下事实推出：
对一个 $A_\infty$ 代数先作 bar 构造、再作 cobar 构造所得的
微分分次代数，同原来的 $A_\infty$ 代数同伦等价；
例如见 \textbf{[Smi00]} 第 4.1 节。

还有函子
\[
[\mathrm{CDGA};\text{同伦}]
\longrightarrow
[\mathrm{DGA};\text{同伦}],
\]
以及显然的函子
\[
\mathfrak{HA}_\infty
=
[\text{滤过 }A_\infty;\text{同伦}]
\longrightarrow
[A_\infty;\text{同伦}],
\qquad
(C,\mathfrak m)
\longmapsto
(\overline C,\overline{\mathfrak m}).
\]
关于记号 $\mathfrak{HA}_\infty$，见第~\ref{subsec:4-3-1} 小节。

\begin{corollary}{7.5.2}\label{cor:7-5-2}
施行 $\otimes\mathbb R$ 后，下列图式交换：

\par\smallskip\noindent\begin{minipage}{\textwidth}
\centering
\resizebox{.99\textwidth}{!}{$\displaystyle
\begin{array}{ccccc}
[\text{相对辛}]
&\longrightarrow&
[\text{滤过 }A_\infty;\text{同伦}]
&\longrightarrow&
[A_\infty;\text{同伦}]
\\[-.1em]
\Big\downarrow
&&&&
{\scriptstyle\cong}\Big\uparrow
\\[-.1em]
[\text{有理同伦}]
&\longrightarrow&
[\mathrm{CDGA};\text{同伦}]
&\longrightarrow&
[\mathrm{DGA};\text{同伦}].
\end{array}$}
\par\smallskip
{\renewcommand{\thefigure}{7.5.1}\refstepcounter{figure}
\normalfont\bfseries 图式~\thefigure\label{diag:7-5-1}}
\end{minipage}\par\smallskip
\end{corollary}

证明由定理~\ref{thm:7-5-1}、\ref{thm:4-1-1}
与~\ref{thm:4-1-2} 立即可得。
